\documentclass[11pt,a4paper]{article}

\usepackage[T1]{fontenc}
\usepackage[utf8]{inputenc}
\usepackage[french]{babel}
\usepackage{lmodern}
\usepackage{geometry}
\usepackage{graphicx}
\usepackage{subcaption}
\usepackage{amsmath,amssymb}
\usepackage{booktabs}
\usepackage{url}
\usepackage{float}
\usepackage{enumitem}

\geometry{margin=2.5cm}

\title{Algorithmique et programmation\\[0.5em]
TD4 : Comparaison d'algorithmes de tri}
\author{Marwane El Bachir Youssef Badr\\L3 MIAGE / CMI / DL}
\date{\today}

\begin{document}

\maketitle
\tableofcontents
\newpage

%=========================================================
\section{Introduction}
%=========================================================

L'objectif de ce projet est de comparer expérimentalement différents
algorithmes de tri étudiés en cours d'algorithmique.  
Plutôt que de se limiter à l'analyse théorique des complexités, on
cherche ici à mesurer les temps d'exécution réels en fonction :

\begin{itemize}
  \item de la \textbf{taille} du tableau à trier~;
  \item de la \textbf{structure initiale} des données (tableau déjà trié,
        trié à l'envers, constant ou aléatoire)~;
  \item de la \textbf{plage de valeurs} des entiers (petites valeurs
        vs grandes valeurs).
\end{itemize}

Les algorithmes étudiés sont :

\begin{itemize}
  \item tri à bulle (trois variantes)~;
  \item tri cocktail~;
  \item tri par sélection~;
  \item tri par insertion~;
  \item tri par comptage (dénombrement)~;
  \item tri fusion~;
  \item tri rapide.
\end{itemize}

Les résultats sont présentés sous la forme de courbes générées à partir
d'un fichier CSV, puis analysés en les confrontant aux complexités
théoriques vues en cours.

%=========================================================
\section{Environnement de test}
%=========================================================

\subsection{Configuration matérielle et logicielle}

Les expériences ont été réalisées sur la machine personnelle suivante
(les valeurs ci-dessous sont indicatives et peuvent être adaptées au
matériel réel)~:

\begin{itemize}
  \item \textbf{Machine} : MacBook Pro (puce Apple Silicon / Intel, 8--16 Go RAM)~;
  \item \textbf{Système d'exploitation} : macOS~;
  \item \textbf{Langage de programmation} : C (standard C99)~;
  \item \textbf{Compilateur} : \texttt{gcc} (via \texttt{clang} sous macOS)~;
  \item \textbf{Options de compilation} : \texttt{-O2} pour activer des
        optimisations classiques.
\end{itemize}

L'utilisation d'un seul cœur est supposée, et aucun autre programme
lourd n'était en cours d'exécution lors des mesures, afin de limiter le
bruit expérimental.

\subsection{Mesure du temps}

Le temps d'exécution a été mesuré en C à l'aide de la fonction
\texttt{clock()} de la bibliothèque standard \texttt{time.h}.  
Pour chaque expérience, on ne mesure que le temps de la \textbf{fonction
de tri}, et non le temps de génération des tableaux de test.

Le principe est le suivant~:

\begin{verbatim}
clock_t debut = clock();
/* appel de l'algorithme de tri */
clock_t fin = clock();
double duree = (double)(fin - debut) / CLOCKS_PER_SEC;
\end{verbatim}

Cette durée (en secondes) est ensuite écrite dans le fichier CSV
\texttt{resultats.csv}, avec les informations associées
(type de tri, taille du tableau, cas, plage de valeurs).

%=========================================================
\section{Algorithmes étudiés}
%=========================================================

Dans cette section, on rappelle brièvement le principe de chaque
algorithme, ainsi que sa complexité théorique en fonction de la taille
du tableau $n$.

\subsection{Tri à bulle et variantes}

Le tri à bulle parcourt le tableau plusieurs fois et échange les
éléments adjacents mal ordonnés.  
La version de base effectue systématiquement $n-1$ passages, indépendamment
de l'ordre des données.

\begin{itemize}
  \item \textbf{Cas moyen et pire} : $\mathcal{O}(n^2)$ comparaisons et échanges~;
  \item \textbf{Meilleur cas} (tableau déjà trié) : toujours $\mathcal{O}(n^2)$
        pour la version basique, car tous les passages sont effectués.
\end{itemize}

Deux optimisations classiques ont été implémentées~:

\begin{itemize}
  \item une version \emph{bulle\_basique\_v2} qui réduit la longueur
        du parcours à chaque itération (les plus grands éléments «~bulle~»
        déjà en fin de tableau)~;
  \item une version \emph{bulle\_optimise} qui s'arrête dès qu'aucun
        échange n'est réalisé pendant un passage, ce qui donne un
        meilleur cas en $\mathcal{O}(n)$ lorsque le tableau est déjà trié.
\end{itemize}

\subsection{Tri cocktail}

Le tri cocktail est une variante du tri à bulle où l'on balaie le
tableau alternativement de gauche à droite puis de droite à gauche.
Cela permet de faire remonter les petits éléments vers le début du
tableau plus rapidement.

La complexité reste globalement $\mathcal{O}(n^2)$, avec cependant un
comportement un peu meilleur sur certains jeux de données grâce aux
deux balayages.

\subsection{Tri par insertion}

Le tri par insertion construit progressivement un tableau trié en
insérant chaque élément à la bonne position dans la partie déjà triée.

\begin{itemize}
  \item \textbf{Pire et cas moyen} : $\mathcal{O}(n^2)$~;
  \item \textbf{Meilleur cas} (tableau déjà trié) : $\mathcal{O}(n)$,
        car chaque nouvel élément est déjà à sa place.
\end{itemize}

Cet algorithme est en pratique très efficace sur de petites tailles ou
sur des données presque triées.

\subsection{Tri par sélection}

Le tri par sélection parcourt le tableau pour trouver le minimum, le
met en première position, puis recommence sur le sous-tableau restant.

Il effectue toujours le même nombre de comparaisons, indépendamment de
l'ordre initial des données~: la complexité est
$\mathcal{O}(n^2)$ dans tous les cas.  
En revanche, le nombre d'échanges est limité (au plus $n-1$).

\subsection{Tri par comptage (dénombrement)}

Le tri par comptage suppose que les entiers à trier appartiennent à un
intervalle connu $[0, K]$.  
On commence par compter, dans un tableau de fréquences, combien de fois
apparaît chaque valeur possible, puis on reconstruit le tableau trié à
partir de ces comptages.

\begin{itemize}
  \item \textbf{Complexité} : $\mathcal{O}(n + K)$~;
  \item très performant lorsque $K$ est petit devant $n$ (par exemple
        valeurs dans $[0,1000]$)~;
  \item beaucoup moins intéressant lorsque $K$ est très grand, car le
        tableau de comptage devient coûteux en temps et en mémoire.
\end{itemize}

\subsection{Tri fusion}

Le tri fusion est un algorithme récursif de type «~diviser pour
régner~».  
Le tableau est coupé en deux, chacun des deux sous-tableaux est trié,
puis les deux résultats sont fusionnés.

\begin{itemize}
  \item \textbf{Complexité}~: $\mathcal{O}(n \log n)$ dans le meilleur,
        le moyen et le pire des cas~;
  \item nécessite un tableau auxiliaire pour la phase de fusion.
\end{itemize}

\subsection{Tri rapide}

Le tri rapide (\emph{quicksort}) choisit un pivot, partitionne le
tableau en éléments plus petits et plus grands, puis trie récursivement
les deux parties.

\begin{itemize}
  \item \textbf{Cas moyen} : $\mathcal{O}(n \log n)$~;
  \item \textbf{Meilleur cas} : également $\mathcal{O}(n \log n)$~;
  \item \textbf{Pire cas} (pivot très mauvais à chaque étape) :
        $\mathcal{O}(n^2)$.
\end{itemize}

En pratique, avec un choix de pivot raisonnable, c'est souvent l'un des
tris les plus rapides sur de grandes tailles.

%=========================================================
\section{Protocole expérimental}
%=========================================================

\subsection{Génération des données}

Les tableaux de test sont générés une fois pour toutes par un programme
C (fichier \texttt{benchmark.c}).  
Pour chaque taille $n$ et pour chaque type de scénario, un fichier texte
est créé.  
On distingue les cas suivants~:

\begin{itemize}
  \item \textbf{Tableau aléatoire}~;
  \item \textbf{Tableau trié croissant}~;
  \item \textbf{Tableau trié décroissant}~;
  \item \textbf{Tableau constant} (toutes les valeurs identiques).
\end{itemize}

Pour chaque cas, deux plages de valeurs ont été utilisées~:

\begin{itemize}
  \item \textbf{Petite plage} : entiers dans un intervalle réduit
        (par exemple $[0,1000]$)~;
  \item \textbf{Grande plage} : entiers dans un intervalle beaucoup
        plus large.
\end{itemize}

Les tailles considérées vont de $n=1000$ jusqu'à $n=49000$, par pas
de 2000 éléments.  
Cela permet d'observer l'évolution du temps d'exécution sur une large
gamme de tailles, tout en restant dans des temps de calcul raisonnables
pour les tris quadratiques.

\subsection{Exécution des algorithmes}

Un second programme C (\texttt{tri.c}) lit les fichiers générés par
\texttt{benchmark.c}, applique chaque algorithme de tri à chaque
instance, mesure le temps d'exécution, puis enregistre les résultats
dans un fichier CSV \texttt{resultats.csv}.

Chaque ligne de ce fichier contient~:

\begin{itemize}
  \item le nom de l'algorithme (\texttt{tri})~;
  \item la taille du tableau (\texttt{taille})~;
  \item le cas testé (\texttt{cas} : aléatoire, croissant,
        décroissant, constant)~;
  \item la plage de valeurs (\texttt{plage} : petite, grande)~;
  \item le temps mesuré en secondes (\texttt{temps}).
\end{itemize}

Au total, le fichier \texttt{resultats.csv} contient plus de 1800
lignes, couvrant toutes les combinaisons de tailles, de cas et
d'algorithmes.

\subsection{Génération des graphiques}

Un script Python (\texttt{plots.py}) lit le fichier
\texttt{resultats.csv} avec \texttt{pandas} et produit des graphiques
au format PNG à l'aide de \texttt{matplotlib}.  
Les fichiers sont enregistrés dans le dossier \texttt{figures/}.

On dispose notamment des figures suivantes~:

\begin{itemize}
  \item pour chaque algorithme de tri~: courbes temps--taille pour la
        plage \emph{petite} et pour la plage \emph{grande}~;
  \item des figures spécifiques pour le tri par comptage (un graphique
        par cas)~;
  \item des figures de comparaison globale des algorithmes pour le cas
        aléatoire.
\end{itemize}

Ces figures sont maintenant intégrées et analysées dans la suite du
rapport.

%=========================================================
\section{Résultats expérimentaux détaillés}
%=========================================================

Dans cette section, on commente les courbes obtenues pour chaque
algorithme.  
Pour chaque tri, on présente deux graphes : un pour la petite plage de
valeurs, un pour la grande.

\subsection{Tri à bulle (trois versions)}

\begin{figure}[H]
  \centering
  \begin{subfigure}{0.48\textwidth}
    \centering
    \includegraphics[width=\linewidth]{figures/tri_bulle_basique_petite.png}
    \caption{Tri à bulle basique (petite plage)}
  \end{subfigure}\hfill
  \begin{subfigure}{0.48\textwidth}
    \centering
    \includegraphics[width=\linewidth]{figures/tri_bulle_basique_grande.png}
    \caption{Tri à bulle basique (grande plage)}
  \end{subfigure}

  \medskip
  \begin{subfigure}{0.48\textwidth}
    \centering
    \includegraphics[width=\linewidth]{figures/tri_bulle_basique_v2_petite.png}
    \caption{Tri à bulle v2 (petite plage)}
  \end{subfigure}\hfill
  \begin{subfigure}{0.48\textwidth}
    \centering
    \includegraphics[width=\linewidth]{figures/tri_bulle_basique_v2_grande.png}
    \caption{Tri à bulle v2 (grande plage)}
  \end{subfigure}

  \medskip
  \begin{subfigure}{0.48\textwidth}
    \centering
    \includegraphics[width=\linewidth]{figures/tri_bulle_optimise_petite.png}
    \caption{Tri à bulle optimisé (petite plage)}
  \end{subfigure}\hfill
  \begin{subfigure}{0.48\textwidth}
    \centering
    \includegraphics[width=\linewidth]{figures/tri_bulle_optimise_grande.png}
    \caption{Tri à bulle optimisé (grande plage)}
  \end{subfigure}
  \caption{Évolution du temps pour les différentes versions du tri à bulle.}
\end{figure}
Pour les trois variantes du tri à bulle, on observe une évolution très caractéristique
de la complexité quadratique. Dans tous les cas, les courbes présentent une croissance
parabolique : lorsque la taille du tableau $n$ est multipliée par deux, le temps
d’exécution est multiplié par un facteur proche de quatre. Ce comportement est
strictement conforme à la complexité théorique en $\mathcal{O}(n^2)$ annoncée en cours.

\paragraph{Version basique}
La version basique effectue systématiquement deux boucles imbriquées, sans optimisation.
Elle traite l’ensemble du tableau même si celui-ci est déjà trié. Les courbes montrent
clairement que :
\begin{itemize}
    \item le cas \textbf{croissant} (meilleur cas théorique) n’apporte aucune amélioration
          car cette version effectue toujours le même nombre de comparaisons ;
    \item le cas \textbf{constant} est également aussi coûteux que les autres cas,
          puisque cette variante ne tient compte ni des égalités ni de l’état de tri ;
    \item le cas \textbf{décroissant} représente le pire cas observé, avec une courbe
          systématiquement supérieure, ce qui correspond au nombre maximal
          d’échanges effectués.
\end{itemize}

\paragraph{Version améliorée (bulle\_basique\_v2)}
Cette variante intègre une optimisation consistant à raccourcir la boucle interne
à mesure que les éléments les plus grands « remontent » en fin de tableau.
Les courbes montrent une légère amélioration :
\begin{itemize}
    \item le cas \textbf{croissant} n’est toujours pas optimisé (pas d’arrêt anticipé),
          mais le nombre de comparaisons est légèrement réduit ;
    \item les cas \textbf{aléatoire} et \textbf{décroissant} montrent un gain modéré,
          car la boucle interne devient un peu plus courte à chaque passage ;
    \item l’ordre relatif des scénarios reste identique à la version basique.
\end{itemize}

\paragraph{Version optimisée (bulle\_optimise)}
C’est la version la plus efficace des trois. Elle intègre un mécanisme d’arrêt anticipé :
si aucune permutation n’est effectuée au cours d’un passage, l’algorithme s’arrête
immédiatement. Les observations expérimentales confirment pleinement l’analyse
théorique :
\begin{itemize}
    \item dans le cas \textbf{croissant}, le temps est presque nul : l’algorithme détecte
          immédiatement que le tableau est déjà trié, ce qui valide le meilleur cas
          en $\mathcal{O}(n)$ ;
    \item dans les cas \textbf{aléatoire} ou \textbf{constant}, les courbes sont très
          proches de celles de bulle\_basique\_v2, car l’optimisation ne s’active que tardivement ;
    \item dans le cas \textbf{décroissant}, on retrouve le pire cas quadratique, comme prévu.
\end{itemize}

\paragraph{Petite plage vs grande plage}
Le passage de la plage « petite » à la plage « grande » ne change pas la forme générale
des courbes. On note cependant que :
\begin{itemize}
    \item l’amplitude des courbes augmente davantage pour les tableaux de grande taille ;
    \item les différences entre les scénarios (croissant, aléatoire, décroissant)
          deviennent plus visibles lorsque $n$ augmente ;
    \item le tri à bulle optimisé conserve une avance nette sur les deux autres versions
          pour les données triées ou presque triées.
\end{itemize}

En résumé, cette première étude confirme rigoureusement que les trois variantes
du tri à bulle restent inutilisables dès que la taille des données dépasse quelques dizaines
de milliers d’éléments, à l’exception de la version optimisée dans le meilleur cas,
où l’arrêt anticipé permet une exécution presque linéaire.

\subsection{Tri cocktail}

\begin{figure}[H]
  \centering
  \begin{subfigure}{0.48\textwidth}
    \centering
    \includegraphics[width=\linewidth]{figures/tri_cocktail_petite.png}
    \caption{Tri cocktail (petite plage)}
  \end{subfigure}\hfill
  \begin{subfigure}{0.48\textwidth}
    \centering
    \includegraphics[width=\linewidth]{figures/tri_cocktail_grande.png}
    \caption{Tri cocktail (grande plage)}
  \end{subfigure}
  \caption{Évolution du temps pour le tri cocktail.}
\end{figure}



Le tri cocktail est une variante du tri à bulle qui effectue des parcours dans les
deux sens : de gauche à droite puis de droite à gauche. Cette amélioration permet
de réduire légèrement le nombre de passes nécessaires lorsque des éléments très grands
ou très petits se situent aux extrémités du tableau.

\paragraph{Analyse des graphes (petite et grande plage)}
Les courbes obtenues présentent la même forme que celles du tri à bulle amélioré :
une croissance quadratique très marquée, confirmant la complexité théorique en
$\mathcal{O}(n^2)$.

On observe que :

\begin{itemize}
    \item les courbes \textbf{croissant} et \textbf{constant} montrent un temps légèrement
    inférieur à la version basique du tri à bulle, mais restent très proches d’une
    croissance en $n^2$ ;
    \item le cas \textbf{décroissant} est le plus coûteux, car les deux parcours
    avant/arrière doivent réorganiser l’ensemble des éléments ;
    \item le cas \textbf{aléatoire} produit des performances intermédiaires ;
    \item la différence entre petite et grande plage est surtout visible dans l’amplitude
    des temps, qui devient plus importante lorsque $n$ dépasse les 20\,000 éléments.
\end{itemize}

\paragraph{Comparaison avec les variantes du tri à bulle}
Comme prévu, le tri cocktail est légèrement plus rapide que le tri à bulle v2 lorsque
de gros éléments sont mal positionnés, mais son comportement reste globalement
équivalent : quadratique, lent et rapidement inutilisable sur de grandes tailles.

En résumé, même avec des passes bidirectionnelles, le tri cocktail n’offre qu’un gain
modéré et reste un tri quadratique à éviter pour des tableaux volumineux.

\subsection{Tri par insertion}

\begin{figure}[H]
  \centering
  \begin{subfigure}{0.48\textwidth}
    \centering
    \includegraphics[width=\linewidth]{figures/tri_insertion_petite.png}
    \caption{Tri insertion (petite plage)}
  \end{subfigure}\hfill
  \begin{subfigure}{0.48\textwidth}
    \centering
    \includegraphics[width=\linewidth]{figures/tri_insertion_grande.png}
    \caption{Tri insertion (grande plage)}
  \end{subfigure}
  \caption{Évolution du temps pour le tri par insertion.}
\end{figure}



Le tri par insertion construit progressivement un tableau trié en insérant chaque
élément à sa position correcte. Sa complexité théorique est en $\mathcal{O}(n^2)$,
mais son meilleur cas est linéaire en $\mathcal{O}(n)$ lorsque le tableau est déjà trié.

\paragraph{Analyse des graphes}
Les résultats expérimentaux montrent exactement ce que prédit la théorie :

\begin{itemize}
    \item dans le cas \textbf{croissant}, le temps d'exécution est extrêmement faible,
    ce qui correspond au meilleur cas linéaire : chaque élément est déjà à la bonne
    position et aucune permutation n’est effectuée ;
    \item dans le cas \textbf{constant}, l’algorithme effectue quelques comparaisons
    mais très peu de déplacements, ce qui explique une courbe légèrement supérieure
    mais encore très faible ;
    \item le cas \textbf{aléatoire} produit une courbe quadratique classique ;
    \item le cas \textbf{décroissant} représente le pire cas théorique, car chaque
    insertion doit décaler l’intégralité du sous-tableau trié.
\end{itemize}

\paragraph{Petite vs grande plage}
Lorsque la plage de valeurs augmente, les courbes conservent la même forme.
La différence se manifeste principalement par des variations aléatoires,
liées aux coûts des comparaisons entre valeurs plus éloignées.

\paragraph{Conclusion}
Le tri par insertion confirme son excellent comportement sur les tableaux
déjà triés ou presque triés. En revanche, son pire cas reste quadratique,
ce qui le rend rapidement inutilisable pour des tableaux supérieurs à 30\,000 éléments.


\subsection{Tri par sélection}

\begin{figure}[H]
  \centering
  \begin{subfigure}{0.48\textwidth}
    \centering
    \includegraphics[width=\linewidth]{figures/tri_selection_petite.png}
    \caption{Tri sélection (petite plage)}
  \end{subfigure}\hfill
  \begin{subfigure}{0.48\textwidth}
    \centering
    \includegraphics[width=\linewidth]{figures/tri_selection_grande.png}
    \caption{Tri sélection (grande plage)}
  \end{subfigure}
  \caption{Évolution du temps pour le tri par sélection.}
\end{figure}



Le tri par sélection cherche à chaque itération le plus petit élément restant
dans le tableau et le place à la bonne position. Sa complexité est rigoureusement
quadratique en $\mathcal{O}(n^2)$, quelle que soit la configuration initiale.

\paragraph{Analyse des courbes}
Les courbes expérimentales illustrent parfaitement cette absence de meilleur cas :

\begin{itemize}
    \item les quatre courbes (croissant, constant, aléatoire, décroissant) sont
    presque parfaitement confondues ;
    \item la différence entre les scénarios est négligeable : le nombre de comparaisons
    est identique pour toutes les configurations ;
    \item sur grande plage comme sur petite, l’évolution est strictement parabolique.
\end{itemize}

\paragraph{Interprétation}
Cette stabilité est logique : le tri par sélection effectue systématiquement
le même nombre d’opérations, que le tableau soit déjà trié ou totalement inversé.

\paragraph{Conclusion}
Bien que simple à implémenter et stable en termes de comportement, le tri par sélection
est l’un des moins performants sur les grands tableaux car il ne bénéficie d’aucun
meilleur cas et reste strictement en $\mathcal{O}(n^2)$.


\subsection{Tri fusion}

\begin{figure}[H]
  \centering
  \begin{subfigure}{0.48\textwidth}
    \centering
    \includegraphics[width=\linewidth]{figures/tri_fusion_petite.png}
    \caption{Tri fusion (petite plage)}
  \end{subfigure}\hfill
  \begin{subfigure}{0.48\textwidth}
    \centering
    \includegraphics[width=\linewidth]{figures/tri_fusion_grande.png}
    \caption{Tri fusion (grande plage)}
  \end{subfigure}
  \caption{Évolution du temps pour le tri fusion.}
\end{figure}



Le tri fusion est l’un des algorithmes les plus performants étudiés ici.
Il repose sur une stratégie diviser-pour-régner et présente une complexité
théorique en $\mathcal{O}(n \log n)$ dans tous les cas.

\paragraph{Analyse des graphes}
Les résultats expérimentaux confirment parfaitement ce comportement :

\begin{itemize}
    \item les quatre cas (croissant, constant, aléatoire, décroissant) produisent
    des courbes quasiment identiques ;
    \item aucune explosion quadratique n’est visible, ce qui contraste fortement avec
    les tris par comparaison quadratiques ;
    \item l’échelle de temps est très faible, même pour les tableaux de 50\,000 éléments.
\end{itemize}

\paragraph{Petite vs grande plage}
L’élargissement de la plage de valeurs n’a pratiquement aucun impact sur les performances :
le tri fusion compare des valeurs, mais la nature des valeurs n’affecte pas le nombre
de comparaisons nécessaires.

\paragraph{Conclusion}
Le tri fusion est l’un des plus fiables et des plus réguliers en termes de performance.
Sa complexité en $\mathcal{O}(n \log n)$ le place parmi les meilleurs tris généraux.


\subsection{Tri rapide}

\begin{figure}[H]
  \centering
  \begin{subfigure}{0.48\textwidth}
    \centering
    \includegraphics[width=\linewidth]{figures/tri_rapide_petite.png}
    \caption{Tri rapide (petite plage)}
  \end{subfigure}\hfill
  \begin{subfigure}{0.48\textwidth}
    \centering
    \includegraphics[width=\linewidth]{figures/tri_rapide_grande.png}
    \caption{Tri rapide (grande plage)}
  \end{subfigure}
  \caption{Évolution du temps pour le tri rapide.}
\end{figure}



Le tri rapide repose également sur une stratégie diviser-pour-régner, mais sa
complexité dépend du choix du pivot. Théoriquement :
\begin{itemize}
    \item meilleur cas : $\mathcal{O}(n \log n)$ ;
    \item cas moyen : $\mathcal{O}(n \log n)$ ;
    \item pire cas : $\mathcal{O}(n^2)$, lorsque le pivot produit des partitions très déséquilibrées.
\end{itemize}

\paragraph{Analyse des courbes}
Les graphes montrent que :

\begin{itemize}
    \item le cas \textbf{aléatoire} produit les meilleures performances ;
    \item le cas \textbf{croissant} ou \textbf{décroissant} peut parfois produire des partitions déséquilibrées :
          on observe certains pics anormaux pour les grandes tailles ;
    \item malgré cela, pour $n < 50\,000$, aucune explosion quadratique n’apparaît,
          ce qui montre que l’implémentation du pivot fonctionne correctement.
\end{itemize}

\paragraph{Impact de la plage}
L’augmentation de la plage de valeurs a un impact très limité, ce qui est normal :
le tri rapide ne dépend pas directement de la valeur des éléments mais de leur ordre relatif.

\paragraph{Conclusion}
Le tri rapide offre un excellent compromis entre simplicité et performance.
Ses rares pics sont cohérents avec la possibilité d’un mauvais pivot,
mais son comportement global reste très proche de celui du tri fusion.


\subsection{Tri par comptage}

Pour le tri par comptage, on dispose de graphes dédiés pour chaque cas.

\begin{figure}[H]
  \centering
  \begin{subfigure}{0.48\textwidth}
    \centering
    \includegraphics[width=\linewidth]{figures/comptage_aleatoire.png}
    \caption{Comptage -- cas aléatoire}
  \end{subfigure}\hfill
  \begin{subfigure}{0.48\textwidth}
    \centering
    \includegraphics[width=\linewidth]{figures/comptage_croissant.png}
    \caption{Comptage -- cas croissant}
  \end{subfigure}

  \medskip
  \begin{subfigure}{0.48\textwidth}
    \centering
    \includegraphics[width=\linewidth]{figures/comptage_decroissant.png}
    \caption{Comptage -- cas décroissant}
  \end{subfigure}\hfill
  \begin{subfigure}{0.48\textwidth}
    \centering
    \includegraphics[width=\linewidth]{figures/comptage_constant.png}
    \caption{Comptage -- cas constant}
  \end{subfigure}
  \caption{Temps du tri par comptage selon le cas testé.}
\end{figure}



Le tri par comptage n’est pas un tri par comparaison. Il est linéaire en
$\mathcal{O}(n + K)$, où $K$ est la taille de la plage des valeurs possibles.

\paragraph{Analyse des graphes}
Les résultats expérimentaux montrent :

\begin{itemize}
    \item une courbe extrêmement faible pour la \textbf{petite plage} (0 à 1000),
          proche d’un temps linéaire ;
    \item une explosion de temps pour la \textbf{grande plage},
          car $K$ augmente fortement et domine la complexité ;
    \item cette différence très marquée entre petite et grande plage confirme
          exactement la formule $\mathcal{O}(n + K)$.
\end{itemize}

\paragraph{Interprétation}
Le tri par comptage est idéal pour des valeurs limitées dans une petite plage.
Mais dès que $K$ devient trop grand, il devient inutilement coûteux en mémoire
et en temps, ce que l’expérience met en évidence de façon spectaculaire.

\paragraph{Conclusion}
Le tri par comptage est le seul algorithme étudié ici dont la performance dépend
directement de la plage des valeurs. Lorsqu’il est applicable, c’est le plus rapide
de tous, mais il reste inutilisable sur une plage trop large.


%=========================================================
\section{Comparaisons globales}
%=========================================================

Pour mieux visualiser les différences entre les familles d'algorithmes,
on dispose de graphes qui superposent plusieurs tris sur le même
scénario.

\begin{figure}[H]
  \centering
  \begin{subfigure}{0.48\textwidth}
    \centering
    \includegraphics[width=\linewidth]{figures/comparaison_aleatoire_petite.png}
    \caption{Comparaison -- cas aléatoire, petite plage}
  \end{subfigure}\hfill
  \begin{subfigure}{0.48\textwidth}
    \centering
    \includegraphics[width=\linewidth]{figures/comparaison_aleatoire_grande.png}
    \caption{Comparaison -- cas aléatoire, grande plage}
  \end{subfigure}
  \caption{Comparaison globale des algorithmes de tri sur le cas aléatoire.}
\end{figure}



Les figures de comparaison montrent trois familles de comportements :

\begin{itemize}
    \item les algorithmes \textbf{quadratiques} ($\mathcal{O}(n^2)$) :
    tri à bulle (toutes versions), tri cocktail, tri insertion, tri sélection ;

    \item les algorithmes \textbf{$\mathcal{O}(n \log n)$} :
    tri rapide, tri fusion ;

    \item le tri \textbf{linéaire} en fonction de $n$ (plus $K$) :
    tri par comptage.
\end{itemize}

À partir d’environ $n = 10\,000$, l’écart devient énorme :
le tri à bulle dépasse plusieurs secondes, alors que le tri rapide et le tri fusion
restent inférieurs à la demi-seconde. Le tri par comptage reste imbattable dans
la petite plage, mais perd tout intérêt lorsque $K$ augmente fortement.

La comparaison petite/grande plage met en évidence que :
\begin{itemize}
    \item la plage a peu d’effet sur les tris quadratiques ;
    \item elle a un impact très fort sur le tri par comptage ;
    \item les tris $n \log n$ restent les plus stables.
\end{itemize}


À partir d'environ $n=10000$, l'écart entre les algorithmes
quadratiques et ceux en $n\log n$ devient énorme~: le temps de tri à
bulle dépasse plusieurs secondes alors que le tri rapide et le tri
fusion restent inférieurs à la demi-seconde.

La comparaison petite/grande plage met également en évidence que~:

\begin{itemize}
  \item la plage de valeurs a peu d'effet sur les algorithmes de
        comparaison (bulle, insertion, sélection, fusion, rapide)~;
  \item en revanche, elle a un impact très fort sur le tri par
        comptage, comme prévu.
\end{itemize}

%=========================================================
\section{Discussion et lien avec la théorie}
%=========================================================

D'un point de vue théorique, on distingue classiquement :

\begin{itemize}
  \item les algorithmes de tri par comparaison en $\mathcal{O}(n^2)$~;
  \item ceux en $\mathcal{O}(n \log n)$~;
  \item et les algorithmes non basés sur la comparaison (comptage,
        radix, etc.) pouvant atteindre $\mathcal{O}(n)$ sous certaines
        hypothèses.
\end{itemize}

Les résultats expérimentaux obtenus sur ce projet sont globalement
\textbf{en accord} avec ces prédictions~:

\begin{itemize}
  \item les tris à bulle, cocktail, insertion et sélection deviennent
        très rapidement impraticables lorsque $n$ dépasse quelques
        dizaines de milliers~;
  \item le tri par insertion reste toutefois intéressant pour de
        petites tailles ou des tableaux presque triés~;
  \item le tri fusion et le tri rapide se comportent beaucoup mieux, et
        sont clairement à privilégier pour des tableaux de grande
        taille~;
  \item le tri par comptage est imbattable lorsque la plage de valeurs
        est petite, mais perd de l'intérêt dès que cette plage
        s'élargit.
\end{itemize}

Quelques écarts peuvent venir de détails d'implémentation (choix du
pivot pour le tri rapide, gestion de la mémoire pour le tri fusion,
optimisations du compilateur, etc.) mais la tendance globale reste très
proche de ce que prédit l'analyse asymptotique.

%=========================================================
\section{Conclusion}
%=========================================================

Ce projet a permis de mettre en pratique l'analyse de complexité des
algorithmes de tri en la confrontant à des mesures expérimentales sur
de véritables données.

Les principaux enseignements sont les suivants~:

\begin{itemize}
  \item la complexité asymptotique ($\mathcal{O}(n^2)$,
        $\mathcal{O}(n \log n)$, $\mathcal{O}(n + K)$) se retrouve très
        clairement dans les courbes de temps~;
  \item pour des tableaux de grande taille, les algorithmes quadratiques
        doivent absolument être évités~;
  \item le tri rapide et le tri fusion offrent un excellent compromis
        entre simplicité d'implémentation et performances~;
  \item le tri par comptage est extrêmement performant lorsque ses
        hypothèses sont satisfaites (plage de valeurs bornée et
        raisonnable).
\end{itemize}

Enfin, l'utilisation combinée de C pour l'expérimentation et de Python
pour l'analyse de données et la génération de graphiques illustre bien
une démarche complète d'évaluation algorithmique, depuis l'implémentation
jusqu'au rapport final.

\end{document}
